\section{TextRef.sty}\label{sec:ref-textref}

Links from the notes into the course textbook. The package holds the
\emph{mechanism} only; which book, where it lives, and which section a
given lecture covers are all course data. Loaded by all three classes,
and inert until a course declares something.

It works in three parts, and the first is the one that carries the weight.

\textbf{Every lecture declares the section it is.} One line at the top of
the file names the coordinate, and the class prints it as a running head
(section~\ref{sec:ref-lecturenotes}). The tie between a lecture and its
section is therefore \emph{structural} rather than something to remember
to write.

If the book is online, that head is also a link --- but the link is
optional, and everything here works without one. A course taught from a
printed text, or from no text at all, uses the same command and never
types a URL.

\textbf{Anything else is a cross-reference.} Pointing at some
\emph{other} section --- usually a prerequisite --- is what
\cs{textbookref} is for, and those resolve through a small manifest of
just the sections a course actually points at.

\textbf{And a book can simply be named.} A syllabus lists its texts; a
handout says ``the algebra text has a fuller treatment''. Neither is
pointing at a section, and both want the book's name and a link to where
it lives. That is the \emph{book registry}, and because naming a book is
not exclusive the way vocabulary is, it is loaded into every document
whole --- see \emph{The book registry} below.

\begin{center}
\begin{tabularx}{\linewidth}{@{}l Y@{}}
\cs{usetext}\texttt{\{apex\}} & in a master's preamble: this document is
  an \emph{apex} document. Loads \file{.course-machinery-local/apex.tex}.\\
\cs{DeclareBook}\texttt{\{name\}\{\dots\}} & a registry entry: what a book
  is called, who publishes it, where it lives. Keyval; see below.\\
\cs{booktitle}\texttt{\{apex\}} & the book's title, set in italic and
  linked to its home page. Starred: unlinked.\\
\cs{bookname}\texttt{\{apex\}} & its short name, for running prose.
  Starred: unlinked.\\
\cs{bookpublisher}\texttt{\{apex\}} & the publisher. Never linked.\\
\cs{bookurl}\texttt{\{apex\}} & the home URL itself, as printed text.\\
\cs{DeclareTextBase}\texttt{\{name\}\{URL\}} & where a book lives; every
  page is appended to it verbatim, so include the trailing slash.
  \cs{DeclareBook}'s \texttt{base} calls this.\\
\cs{SetDefaultText}\texttt{\{name\}} & the book assumed when a lecture
  does not say. Default: unnamed.\\
\cs{SetTextBase}\texttt{\{URL\}} & shorthand for a course with one book:
  the base for whichever text is the default. (Reads the default name
  \emph{when called}, so if \cs{SetDefaultText} is used too, it must come
  first; \cs{DeclareTextBase} sidesteps the ordering entirely.)\\
\cs{booksection}\texttt{\{2.5\}\{Title\}} & this lecture's own
  coordinate: its title and its running head. No URL needed.\\
\cs{booksection}\texttt{[link=page]\dots} & the same, linked:
  \texttt{link} is a page \emph{appended to the base}, not a full URL.\\
\cs{booksection}\texttt{[text=name]\dots} & the same, naming a book other
  than the default. Combines with \texttt{link}.\\
\cs{lecturetitle}\texttt{\{Title\}} & for a lecture with no counterpart
  in the book: title only, no coordinate, no link.\\
\cs{DeclareSection}\texttt{\{5.2\}\{page\}} & one cross-reference entry:
  section number $\to$ its page. Last declaration wins.\\
\cs{DeclareSection}\texttt{[name]\{3.2\}\{page\}} & the same, for a
  \emph{named} book rather than the document's own.\\
\cs{textbookref}\texttt{\{5.2\}\{Table 5.2.3\}} & prints
  ``Table 5.2.3'', linked to \S5.2.\\
\cs{textbookref*}\texttt{\{5.2.6\}\{Example\}} & starred: appends the key,
  so this prints ``Example~5.2.6''.\\
\cs{textbookref}\texttt{[name]\{3.2\}\dots} & resolved in a named book,
  against that book's base. Star comes first: \cs{textbookref*}\texttt{[name]}.\\
\end{tabularx}
\end{center}

\subsection*{The lecture's own coordinate}

\begin{manexsource}
\documentclass{LectureNotes}
\begin{document}
\booksection{2.5}{The Chain Rule}
\end{manexsource}

That is the whole of it for a course with no online text: the head reads
``2.5: The Chain Rule'', and no base, no text layer and no URL are
involved anywhere. It replaces the \cs{section*} a lecture used to open
with --- the head \emph{is} the title, so there is one place the title is
written and no way for head and body to drift apart.

A course whose book is online adds the page, and nothing else changes:

\begin{manexsource}
\documentclass{LectureNotes}
\begin{document}
\booksection[link=sec_chainrule.html]{2.5}{The Chain Rule}
\end{manexsource}

The base comes from the course's text layer, which is loaded for you --- see
\emph{Loading the text layer} below. Note there is still no \cs{input}.

\textbf{The page lives in the lecture, not in a shared list.} That looks
like the thing this package was built to avoid --- a URL per file, each
rotting on its own schedule --- but it is not. It is one page per file,
in a fixed position, and it can only rot when the book restructures,
which renumbers the section too, so you are editing that line anyway. The
failure that genuinely wants central handling is a publisher moving
domains, and that is still one edit: the base is declared once and only
the page is local.

\textbf{Options, not more arguments.} \texttt{link} and \texttt{text} are
keyval options for the same reason \cs{importproblem} takes
\texttt{points=} and \texttt{name=} that way --- and because it is what
lets the link be genuinely optional, instead of an empty \texttt{\{\}} in
every lecture of every course that has no URL to give.

\textbf{The number is the book's, not the course's.} A course that
teaches a topic earlier than its text does still cites where the material
\emph{lives} --- a lecture eighth in the term can perfectly well carry
\cs{booksection}\texttt{\{5.1\}}. Sequence and text coordinate are
independent; the filename carries one, this line the other.

\subsection*{Loading the text layer}

A course's text data lives in \file{.course-machinery-local/}, the course's
tier-two machinery folder (section~\ref{sec:extending}), and \textbf{no
master \cs{input}s it by path.} \file{.latexmkrc} puts that one folder on
\env{TEXINPUTS}, so the file is found \emph{by name} from any depth. That
is what lets a lecture sit anywhere in the tree: there is no \texttt{../..}
in a master to keep true to the folder layout, and moving a file deeper or
shallower breaks nothing. Only what is deposited in that folder is exposed
this way, not the rest of the course.

Two things load from there, and they load differently.
\file{books.tex}, the registry, is loaded \textbf{always and whole}, into
every document. The \emph{text layer} is the one a master picks:

\begin{itemize}
  \item \textbf{Nothing} --- \file{textbook.tex} is loaded if it exists.
    The ordinary case: a course with one book writes no \cs{usetext}
    anywhere.
  \item \cs{usetext}\texttt{\{apex\}} --- loads \file{apex.tex} instead.
    For a course drawing chapters from more than one book.
  \item Neither, and no \file{textbook.tex} --- nothing loads, and nothing
    breaks. A course with no online text needs no text layer at all, and
    may have no \file{.course-machinery-local/} folder either if nothing
    else of its own is shared across files; \cs{booksection} still sets
    the head.
\end{itemize}

\file{books.tex} is optional on the same terms: absent, nothing loads and
nothing breaks. It is loaded from \emph{both} the \cs{usetext} call and
the end-of-preamble hook, guarded so the second is a no-op --- which is
what lets \cs{booktitle} and \cs{booksection} be used in a preamble as
well as a body, without an ordering rule to remember.

\textbf{The presence of \file{textbook.tex} is what declares whether the
course has a default book}, so there is no separate switch. A course with
two books simply omits it, which makes naming the book in each lecture
mandatory rather than merely possible --- worth having, because a lecture
then carries its book with it, and ``which lectures use the other book?''
is a positive grep rather than a search for absence. If you would rather
have a default anyway, \file{textbook.tex} containing nothing but
\cs{usetext}\texttt{\{apex\}} gives one that a lecture can still override.

\textbf{Why a document names its book, rather than each reference.} Two
kinds of fact attach to a textbook, and only one of them is exclusive.
A book's \emph{identity} --- its title, its publisher, where it lives,
which page a section sits on --- is not: \cs{booksection}\texttt{[text=\dots]}
has always resolved a base per call, and a syllabus needs to name three
books at once. But the notes noun and the \cs{NewTextResult} vocabulary
are document-wide --- a document cannot call worked items ``Example'' and
``Exercise'' at once, nor declare one result environment twice. So the
document, not the reference, is the thing that has to choose, and
\cs{usetext} chooses \emph{only} that. That granularity is also what makes
the duplicate-vocabulary error unreachable: two books may both name a
result, because they are never loaded together.

Identity, being the half that is not exclusive, lives in the registry
instead.

\subsection*{The book registry}

\begin{manexsource}
\DeclareBook{algtrig}{%
  short     = Algebra \& Trigonometry,
  title     = Algebra and Trigonometry 2e,
  publisher = OpenStax,
  home      = https://openstax.org/details/books/algebra-and-trigonometry-2e,
  base      = https://openstax.org/books/algebra-and-trigonometry-2e/,
}
\end{manexsource}

\file{.course-machinery-local/books.tex} holds one such block per book,
and every document gets all of them. Only \texttt{base} is
machinery-facing --- it is what \cs{booksection} and \cs{textbookref}
append pages to, and it simply calls \cs{DeclareTextBase}. The other four
exist to be printed:

\begin{center}
\begin{tabularx}{\linewidth}{@{}l Y@{}}
\texttt{short} & what to call it in running prose (``APEX'').\\
\texttt{title} & the book's actual title.\\
\texttt{publisher} & who publishes it, where that is worth saying.\\
\texttt{home} & the landing page a reader is sent to. Defaults to
  \texttt{base}, which is right whenever a book's web edition is also its
  front door; give both when the publisher keeps a separate catalogue
  page.\\
\texttt{base} & what pages are appended to. Trailing slash included.\\
\end{tabularx}
\end{center}

\texttt{short} and \texttt{title} each default to the other, and to the
book's \emph{key} if neither is given --- so a bare entry prints
\texttt{algtrig} rather than nothing, which is a visible placeholder
instead of a silent gap.

\textbf{The registry is what gives a third book somewhere to live.} A
course often points at a text it never teaches from --- the prerequisite
algebra book behind a calculus course. That book has identity and no
vocabulary at all, so no document is ever ``an \emph{algtrig}
document'' and no \cs{usetext} would ever load a file for it. Before the
registry, every reference to such a book was a hand-typed \cs{href} with
the publisher's domain baked into it, which is the exact rot this package
exists to prevent.

\textbf{Keyval caveat.} A value containing a bare comma splits, and the
error names the fragment: \texttt{`Volume 1' undefined in families
`declarebook'}. Brace it --- \texttt{title = \{Calculus, Volume 1\}}.
\texttt{\#} and \texttt{\%} still need escaping, as they do in a slug.

\subsection*{Three decisions that shape it}

\textbf{Section granularity, never finer.} Table 5.2.3, Key Idea 5.2.1
and Example 5.2.6 all resolve to the top of \S5.2, and the student
scrolls --- which is a feature, not a concession: it puts the section's
context in front of them. Page URLs are also the stable part of a web
text, where in-page anchors are exactly what rots, and it survives a
re-edition that renumbers items \emph{within} a section. A full
coordinate is accepted anyway --- \cs{textbookref}\texttt{\{5.2.3\}}
truncates to \texttt{5.2} on a miss --- so you can type what you read off
the page.

\textbf{Graceful degradation is the point.} The text always prints. It
becomes a link only if the base (and, for \cs{textbookref}, the key)
resolves \emph{and} hyperref is loaded. An undeclared book, a typo'd key,
or an Assessment master (which deliberately loads no hyperref --- a link
on paper is a dead end) all demote the reference to correct plain text.
Nothing breaks, and nothing is silently wrong.

\textbf{Misses are reported.} Degrading silently would let a whole rotted
set of links pass unnoticed, since the document still compiles and still
reads correctly. So each unresolved key or book gets one \cs{PackageInfo}
in the \file{.log} (deduplicated --- a section cited twenty times logs
once), and the end of the run warns with the full list. Both kinds of
failure share one list, so a build is never half-warned:

\begin{manexsource}
Package TextRef Warning: Textbook references printed WITHOUT links
(no matching \DeclareSection or \DeclareTextBase):
(TextRef)                text `nosuchbook', 9.9.
\end{manexsource}

\subsection*{The layering rule}

\cs{textbookref} is a \emph{text-tied} citizen, and the model only holds
if it stays where the text is known:

\begin{itemize}
  \item \textbf{Yes:} lecture bodies, and box titles written at the call
    site.
  \item \textbf{No:} inside \file{CourseBoxes.sty} or
    \file{Exercises.sty} -- that would couple text-agnostic machinery to
    a particular book.
  \item \textbf{No:} inside a \file{problems/} or \file{facts/} file.
    Those are reused across courses and texts; a baked-in reference would
    point at the wrong book the moment the file is borrowed.
  \item \textbf{No:} an import key. \env{name=} is a short \emph{content
    label} -- it titles the box and becomes the contents-list entry, and a
    link nested inside a contents entry's own link is not something to
    invite. Provenance is \env{source}/\env{number} in the file's metadata
    block, and the displayed \env{number} already says where it came from.
\end{itemize}

The reference therefore goes in the lecture \emph{around} the import --
prose the author writes where the text is known:

\begin{manexsource}
The rule below is stated as \textbookref{2.1}{APEX 2.1}:

\importfact[kind={theorem}, name={Power Rule}, number={2.1}]%
  {../facts/}{power-rule.tex}
\end{manexsource}

\subsection*{A course drawn from two books}

An opening review chapter out of one text and the rest out of another is
a normal situation, and it needs no special handling. Declare both bases,
and let each lecture name the book it means:

\begin{manexsource}
\DeclareTextBase{apex}{https://example.org/apex/}
\DeclareTextBase{openstax}{https://example.org/openstax/}
\SetDefaultText{apex}
\end{manexsource}

\begin{manexsource}
\booksection[text=openstax, link=review-of-functions]{1.1}{Review of Functions}
\end{manexsource}

Because each lecture carries its own coordinate, two books whose section
numbers collide --- and they will, every book has a \S1.1 --- cannot be
confused for one another. The book is part of what the lecture declares.

\textbf{An unqualified cross-reference belongs to the text; a qualified
one belongs to the course.} This is the easy thing to get wrong. A lecture
that says \cs{usetext}\texttt{\{openstax\}} never loads the other book's
file, so that file's plain \cs{DeclareSection} entries are simply
unavailable to it --- an unqualified declaration means ``a section of the
book \emph{this document} is built on'', which is the only thing a file
loaded by \cs{usetext} can honestly say. Each text file therefore carries
the sections \emph{its own} lectures point at.

Naming the book lifts exactly that restriction:

\begin{manexsource}
\DeclareSection[algtrig]{3.2}{pages/3-2-domain-and-range}
...then in any lecture, whatever it is built on:
\textbookref[algtrig]{3.2}{\S3.2, Domain and Range}
\end{manexsource}

Those go in \file{books.tex}, beside the \cs{DeclareBook} they qualify,
because the documents that need them are built on some other book. The two
namespaces \textbf{do not cross}: \cs{DeclareSection}\texttt{\{1.3\}} in
\file{apex.tex} is not reachable as
\cs{textbookref}\texttt{[apex]\{1.3\}}, since resolving the empty name as
``apex'' would be a guess about a file any document may have loaded. A
section wanted both ways is declared twice.

One thing still cannot be shared between two books in a single document:
\emph{vocabulary}. The box factory
(section~\ref{sec:ref-courseboxes}) ends in a fresh environment
definition, which errors on a name already taken --- so two books that
both name a result cannot both declare it in one master. That is a loud
failure: it stops the build and tells you. In practice it does not arise,
since students read one book and the second is merely somewhere you
point.

\subsection*{Why no URL registry}

This package writes nothing at compile time for anything else to read
back, and deliberately so. The \emph{sources} are the record, in three
places: the \cs{booksection} line in each lecture; the base and
cross-reference lines in the course's text layer; and every \cs{demo}
call (section~\ref{sec:ref-titleblock}). Same discipline as ProblemMeta and
FactMeta --- ``what does this course contain?'' must be answerable
without running \LaTeX.

That is also the reason for one rule of style: \textbf{keep each
declaration on one line}. It is what makes the answer a \texttt{grep} ---
every link in a course can be listed, counted or audited without
compiling it, and feeding them all to \texttt{curl} to find what rotted
over the summer is then a short script rather than a project. The rule
pays for itself on readability alone, so it holds whether or not such a
script is ever written.

\subsection*{For package authors}

Nothing in this package typesets. \cs{booksection} /
\cs{lecturetitle} only \emph{record}, into public read-only accessors:
\cs{BookSectionHead} (the assembled head --- number, separator, title,
linked if a URL resolved and hyperref is present), and the raw pieces
\cs{BookSectionNumber}, \cs{BookSectionTitle}, \cs{BookSectionText} and
\cs{BookSectionSlug}. A class that wants the running head reads
\cs{BookSectionHead} at shipout --- LectureNotes puts it in
\cs{CourseRunningHead} (TitleBlock) and keys its \cs{headrule} off its
emptiness, so a lecture that declares nothing gets no head and no rule
--- and a class that wants no head simply never reads it: the same
mechanism/presentation split as everywhere else
(section~\ref{sec:hooks}). \cs{BookSectionHead} is
\cs{providecommand}-stubbed empty in TitleBlock, so reading it never
requires this package to be loaded. The one knob between the pieces is
\cs{BookSectionSeparator} (section~\ref{sec:ref-knobs}).
